\section{造价模型探索：原理、数据准入与完整试验}\label{app:cost-neural}

\subsection{论文提供的方法与本项目的预测时点}

用户提供的王宝龙等研究以广东200个EPC改造工程为样本：广州50个用于训练和调参，深圳、佛山、东莞150个用于外部测试。论文报告广州MAPE为0.31\%、$R^2=0.9968$，其他三市的表现见原文实验；这些数值归属于论文研究。本项目已完成中标价格探索训练与七指标合成示例，真实EPC结算模型及浙江测试仍待样本。\cite{psobp2026}

论文的目标是EPC竣工单方成本，输入包括招标阶段分部分项实体费用。因此接入分为两个时点：前期策划沿用带来源的工程量、单价与同口径历史参考组；取得控制价清单及工程条件后，再研究竣工成本预测。工程范围分别登记施工与EPC；中标、合同、结算分别保存。当研究目标改为中标或合同单方价格时，另行建立对应标签与验证结果。

\subsection{七个输入如何采集}

\begin{longtable}{@{}L{12mm}L{29mm}L{58mm}L{55mm}@{}}
\toprule\textbf{输入}&\textbf{单位与含义}&\textbf{应取得的原始资料}&\textbf{首版148条审计的情况}\\
\midrule\endfirsthead
\toprule\textbf{输入}&\textbf{单位与含义}&\textbf{应取得的原始资料}&\textbf{首版148条审计的情况}\\
\midrule\endhead
X1&改造内容，元/平方米&分部分项招标控制实体费用与同范围建筑面积&有107条控制价总额，尚无独立实体费用字段\\
X2&改造深度，覆盖指数&六类子系统覆盖度及当地估算指标权重&有公开改造内容文本，覆盖度和权重待采集\\
X3&施工条件，指数&最小道路宽度、临时堆放及理论临时面积&三个原始字段均缺失\\
X4&使用年限，年&建设年份与预测基准年&建设年代原文1条，需逐项明确年份\\
X5&建筑面积，平方米&与造价分子同工程范围的面积说明&61条有建筑面积，需核对金额与面积匹配范围\\
X6&其他费用，元/平方米&其他费用明细与同范围建筑面积&没有独立其他费用字段\\
X7&工期，天&预测时点的招标、合同工期及出处&招标或中标工期115条，合同期限18条，分别保留\\
\bottomrule\end{longtable}

2026年10月14日首版工作簿审计共148条工程记录，早于正文采用的1277条补全库。其中浙江52、上海50、广东46；合同金额22、中标金额123、控制价107、建筑面积61、竣工结算金额0条。合同金额记录中，EPC14条、施工8条。该统计来自本次只读脚本，不将表内“已核验”标注当作本次重新核验148份公告。原件哈希及字段统计保存在\path{data/imports/cost-neural-dataset-audit-20261014.json}，执行入口为\path{scripts/audit-cost-neural-inputs.py}。

当前准入结果为缺少竣工结算标签和关键输入。先按项目及标段关联官方公告、控制价清单、合同、结算和面积原文，补齐价格年、税费及施工范围，形成可追溯的项目行；AI模拟数据用于界面与计算流程演示，真实造价预测效果使用独立历史项目验证。

\subsection{广东参数怎样作为浙江的研究起点}

需要分别记录三类参数。第一类为算法搜索设置：论文采用30个粒子、100次迭代、学习因子1.7与1.5、惯性权重0.9至0.4、五折验证，搜索节点数、学习率和正则化系数；这些可作为首轮实验配置。第二类为神经网络学习得到的权重、偏置及归一化器：现有PDF没有提供完整可加载模型和原始200条训练表，需要作者模型文件或依据统一数据重新训练。第三类为工程指数中的子系统权重与施工条件组合系数：论文改造深度引用广东估算指标，浙江优先采用当地依据；暂用广东系数时登记出处，并在浙江样本上比较敏感性。SHAP贡献排序单独作为解释结果保存。\cite{psobp2026}

论文原文说明样本来自广东公共资源交易平台。作者单位包括吉林省建苑设计集团有限公司和长春工程学院；单位信息不作为企业内部样本来源的证据。公开公告能够提供项目规模、价格阶段和工期，分部分项、现场条件与结算仍需逐份原件取得。

\subsection{浙江验证与真实数据回流}

训练前固定工程范围、预测时点、标签定义、单位和数据版本，按项目分组隔离同一工程的多个公告。浙江留出项目按城市或时间设计独立测试，训练部分内部进行五折PSO调参，测试部分不参与搜索和归一化器拟合；如采用交叉验证估计最终性能，使用外层评价与内层调参。调参与性能估计使用同一组验证结果会产生乐观偏差，嵌套交叉验证的方法依据见官方文档。\cite{nestedcv}

同一测试集比较同口径单方中位数、岭回归、普通BP与PSO--BP，报告样本数、MAE、RMSE、MAPE及$R^2$；按工程范围、城市和价格年查看误差。数据充分后设置独立校准集估计预测区间。工程人员复核异常项目的面积、价格阶段和变更原因，保留修正前后值与原件版本。新的结算资料审核后进入下一版训练库，每次模型更新保存数据划分、随机种子、配置、模型及测试报告，依据浙江外部测试比较结果选择版本。

造价探索沿用现有解析、证据定位、价格阶段隔离和静态核算。数据准入、中标价试验与合成示例分别记录；新补全数据按原件继续完善特征，浙江预测效果使用同范围浙江项目验证。

\subsection{已执行的公告数据试验与七指标示例}

为测试现有数据的可用性，选取37个具有施工中标金额及受益建筑面积的工作簿项目，目标为中标金额除以受益建筑面积；施工加运营、施工合同与EPC合同分别保留。37个项目包含屋面8、外墙2、卫生设施1、综合整治9、范围尚待明确17个。上海28个分为训练20、开发4、另行评估4；广东9个作为锁定外部测试。仅使用建筑面积对数与原文登记的范围类别，未输入中标金额、单方价或其衍生比值。由于数据由AI采集，原公告逐条复核仍需完成。

另生成480条七指标合成数据，输入范围及成本生成公式另存于运行目录。合成目标由实体费用、改造覆盖、施工条件、年限、面积、其他费用、工期的明确组合加8\%随机扰动生成；300条训练，开发、另行评估和锁定测试各60条。全部模拟标识随模型和导出保存。

两组均采用两个tanh隐藏层及log价格输出。BP以Adam更新权重，外层PSO采用5个粒子、5代搜索，优化隐藏层宽度、学习率和L2；每次拟合500步，种子固定20261014，缩放参数只按训练行拟合。CPU适配复用现有NumPy，不直接使用论文MATLAB trainlm及作者网络权重。\cite{adam} 开发集用于选参，锁定测试一次评估后不反馈搜索。模型JSON重载后预测与保存结果核对一致。

\begin{longtable}{@{}L{36mm}L{26mm}L{29mm}L{29mm}L{26mm}@{}}
\toprule\textbf{锁定数据/模型}&\textbf{样本数}&\textbf{MAE（元/平方米）}&\textbf{MAPE（\%）}&\textbf{$R^2$}\\
\midrule\endfirsthead
\toprule\textbf{锁定数据/模型}&\textbf{样本数}&\textbf{MAE（元/平方米）}&\textbf{MAPE（\%）}&\textbf{$R^2$}\\
\midrule\endhead
广东/中位数基线&9&248.52&179.30&$-2.7546$\\
广东/岭回归&9&253.36&146.49&$-3.6140$\\
广东/BP&9&735.33&321.47&$-85.6328$\\
广东/PSO--BP&9&214.38&128.77&$-1.6555$\\
合成/中位数基线&60&395.36&41.78&$-0.0013$\\
合成/岭回归&60&140.84&11.44&0.8717\\
合成/BP&60&85.87&6.62&0.9406\\
合成/PSO--BP&60&78.09&6.28&0.9452\\
\bottomrule\end{longtable}

公告试验的外部表现较弱。广东9个均为训练未见的综合整治类别，地区、价格年及范围差异同时存在；负$R^2$及128.77\%的百分比误差显示，该模型还不能可靠迁移。浙江的建筑面积样本为0，本次没有浙江独立验证。合成模型的6.28\%只说明对生成机制的学习效果，供参数调整、权重读回和前向推理演示。

网页默认使用七指标合成模型；可切换公告试验，查看相同测试集上四种方法的完整结果。调整输入立即重算单方价及按面积折算的总金额，导出保留目标阶段、模型版本、数据类别和当前参数，真实工程预算保持原计算口径。执行记录位于\path{research/training/cost-neural-prototype/}，模型入口为\path{scripts/train-cost-neural-prototype.py}。本轮运行约3.05秒，峰值内存44.22MiB，不占用GPU；首次Windows内存采样失败及修正后的子实验合同分别保存。
